\section{Abstract Factory}
\label{sec:pat-abstract-factory}

\problemstatement{Provide an interface for creating \emph{families} of
related objects without specifying their concrete classes, so an entire
family can be swapped in one move -- for example a whole ``dark theme'' set
of UI widgets versus a ``light theme'' set.}

\begin{patternbox}[When You Reach for It]
The trigger is \emph{``several related objects that must be used
together and vary as a set.''} A Windows button must pair with a Windows
checkbox, not a Mac one. When products come in matched families and you
must guarantee consistency within a family, reach for Abstract Factory.
\end{patternbox}

\subsection*{Understanding the pattern}

Where Factory Method creates \emph{one} product, Abstract Factory creates a
\emph{suite}. You define an abstract factory interface with a creator
method per product type (\texttt{createButton}, \texttt{createCheckbox}).
Each concrete factory (\texttt{WindowsFactory}, \texttt{MacFactory})
produces the matching family. Client code holds an abstract factory and
asks it for products, never knowing or mixing platforms.

\begin{ideabox}[A Factory of Factories, One per Family]
Group the creators for a coherent family behind a single factory
interface. Selecting a concrete factory once (at start-up) fixes the whole
family, structurally preventing a mismatch like a Mac button beside a
Windows checkbox.
\end{ideabox}

\subsection*{C++ implementation}

\begin{lstlisting}
#include <bits/stdc++.h>
using namespace std;

struct Button   { virtual void paint() = 0; virtual ~Button()   = default; };
struct Checkbox { virtual void paint() = 0; virtual ~Checkbox() = default; };

struct WinButton   : Button   { void paint() override { cout << "Win button\n"; } };
struct WinCheckbox : Checkbox { void paint() override { cout << "Win checkbox\n"; } };
struct MacButton   : Button   { void paint() override { cout << "Mac button\n"; } };
struct MacCheckbox : Checkbox { void paint() override { cout << "Mac checkbox\n"; } };

struct GUIFactory {                            // abstract factory
    virtual unique_ptr<Button>   createButton()   = 0;
    virtual unique_ptr<Checkbox> createCheckbox() = 0;
    virtual ~GUIFactory() = default;
};
struct WinFactory : GUIFactory {
    unique_ptr<Button>   createButton()   override { return make_unique<WinButton>(); }
    unique_ptr<Checkbox> createCheckbox() override { return make_unique<WinCheckbox>(); }
};
struct MacFactory : GUIFactory {
    unique_ptr<Button>   createButton()   override { return make_unique<MacButton>(); }
    unique_ptr<Checkbox> createCheckbox() override { return make_unique<MacCheckbox>(); }
};

void renderUI(GUIFactory& f) {                 // platform-agnostic
    f.createButton()->paint();
    f.createCheckbox()->paint();
}

int main() {
    WinFactory win; MacFactory mac;
    renderUI(win);
    renderUI(mac);                             // whole family swapped
    return 0;
}
\end{lstlisting}

\begin{complexitybox}[Trade-offs]
\textbf{Pros.} Guarantees products from one family are used together;
swapping families is a one-line change; isolates concrete classes.
\textbf{Cons.} Adding a \emph{new product type} means changing the factory
interface and every concrete factory -- a real cost when the product set
grows. Best when families change often but the set of product types is
stable.
\end{complexitybox}

\begin{takeawaybox}
Abstract Factory is Factory Method scaled to families: one factory
interface with several creators, one concrete factory per family. It buys
family consistency at the price of rigidity when new product \emph{types}
appear. Recognise it whenever ``these objects must all match'' is a
requirement.
\end{takeawaybox}
